\chapter{Optimisations Tried in the Kernel Lab}\label{r03:ch:levers}

\emph{Every optimisation that was tried, measured on one kernel in the kernel lab, with its
effect and the decision taken.}

All figures in this chapter are for $64^3$ cells, one core, GCC~14 with \code{-O3}, unless a
table says otherwise. Each table comes from one session, and speedups are only computed within
a session, because the board's absolute speed varies by a few percent between sessions
(\cref{r03:sec:bg-measure}). Tables of AA kernels say whether they show the fastest step or the
cycle average (\cref{r03:sec:lab}). ``Cycles per cell'' are always averages over all steps.

\begin{table}[htbp]
  \centering\small
  \caption{Optimisations tried in the lab, and the decisions.}\label{r03:tab:levers}
  \begin{tabularx}{\linewidth}{@{}p{4.4cm}lX@{}}
    \toprule
    \textbf{Optimisation} & \textbf{Section} & \textbf{Outcome}\\
    \midrule
    software prefetch & \ref{r03:sec:lever-pf} & f32: 653 $\rightarrow$ 311 cycles per cell;
      \textbf{adopted} as a \ps{} option\\
    kernel without the density/velocity output & \ref{r03:sec:lever-pf} & 653 $\rightarrow$ 559 on the
      original kernel; implemented in sweepgen and the application; \textbf{pending}\\
    direction split (hand kernel) & \ref{r03:sec:lever-split} & 6.98 against 5.2~MLUP/s for the
      generated kernel with prefetch; led to the next two rows\\
    AA streaming & \ref{r03:sec:lever-aa} & hand kernel: 180 cycles per cell; generated with
      prefetch: \sx{1.60} (f32), \sx{1.58} (f64) over pull with prefetch; \textbf{adopted}
      in sweepgen and waLBerla\\
    register grouping LMUL 2 and 4 & \ref{r03:sec:lever-sched} & slower (216 against 163 cycles
      per cell); \textbf{rejected}\\
    X60 scheduling model (Clang) & \ref{r03:sec:lever-sched} & +24\,\% on the hand kernel, but GCC
      stays ahead on generated code; \textbf{adopted} as a build flag for Clang\\
    .vf operands and FMA & \ref{r03:sec:lever-selector} & f32 AA with prefetch: 234 $\rightarrow$
      192 cycles per cell; \textbf{adopted} in \ps{}\\
    row pointers, constants in registers & \ref{r03:sec:lever-rp} & lab: fewer instructions, same
      cycles; application: +2.5\,\% (f32), +2.2\,\% (f64); \textbf{adopted} in \ps{}\\
    alignment to 32/64~bytes, huge pages & \ref{r03:sec:lever-layout} & 16~bytes (waLBerla's
      layout) is best for the AA kernels; \textbf{rejected}\\
    padding between the PDF slices & \ref{r03:sec:lever-layout} & \textbf{pending}\\
    direction split, generated from lbmpy & \ref{r03:sec:lever-gensplit} & correct under QEMU;
      \textbf{pending}\\
    \bottomrule
  \end{tabularx}
\end{table}

\Cref{r03:fig:ladder} shows the lab kernels of this chapter on one scale.

\begin{figure}[htbp]
  \centering
  % Lab kernels, f32, kernel alone: cycles per cell (cycle counters; lower is better). Sessions differ by a
% few percent. Sources: lab/results.tsv, lab3.log, lab4.log (appendix D). Target: 8x the application's plain scalar build
% (1.398 MLUP/s, run B) = 1600 / (8 * 1.398) = 143 cycles per cell.
\begin{tikzpicture}
  \begin{axis}[
      xbar, bar width=7pt, bar shift=0pt,
      width=0.66\linewidth, height=7.4cm,
      xmin=0, xmax=1180,
      xlabel={cycles per cell, f32, kernel alone (lower is better)},
      symbolic y coords={hcl,hgcc,gaan,gaao,split,pf,expl,scal},
      ytick={hcl,hgcc,gaan,gaao,split,pf,expl,scal},
      yticklabels={{hand AA + split, Clang with X60 model},{hand AA + split, GCC},
        {generated AA + prefetch, new instr. selection},{generated AA + prefetch, original selection},
        {hand direction split (pull)},{generated pull + prefetch 128 B},
        {generated pull (before this work)},{scalar with CSE}},
      yticklabel style={font=\scriptsize,align=right},
      tick label style={font=\scriptsize}, label style={font=\footnotesize},
      enlarge y limits=0.08,
      nodes near coords, nodes near coords style={font=\tiny,anchor=west},
      xmajorgrids, grid style={psgray!20},
      clip=false,
      extra x ticks={143}, extra x tick labels={},
      extra x tick style={grid=major,major grid style={r03rejected,dashed,thick}},
    ]
    \addplot[fill=psgray!45,draw=psgray!80!black] coordinates {(1035,scal) (653,expl)};
    \addplot[fill=r03landed!55,draw=r03landed!80!black] coordinates {(311,pf) (192.1,gaan)};
    \addplot[fill=psblue!35,draw=psblue!70!black] coordinates {(237,split) (163.9,hgcc) (147.8,hcl)};
    \addplot[fill=r03landed!25,draw=r03landed!70!black] coordinates {(234,gaao)};
    \node[font=\scriptsize,text=r03rejected!80!black,anchor=south west] at (rel axis cs:0.1212,1)
      {\sx{8} the plain scalar build: 143 cycles per cell};
  \end{axis}
\end{tikzpicture}

  \caption[Lab kernels on one scale]{Lab kernels, f32: cycles per cell of the kernel alone
    (lower is better), averaged over all steps. Grey: the kernels before this work; teal:
    adopted; light teal: the generated AA kernel before the instruction-selection change; blue:
    hand kernels. The values come from several sessions, which differ by a few percent. Dashed:
    the cycles per cell that \sx{8} the plain scalar application would need.}\label{r03:fig:ladder}
\end{figure}

\section{Software prefetch}\label{r03:sec:lever-pf}

\code{edits.py} inserts one prefetch per distinct field access of the cell loop, \code{D}
bytes ahead of the address the loop accesses in that iteration: \code{prefetch.r} for the 19
PDF reads, \code{prefetch.w} for the 23 stores. The kernels are compiled with
\code{-march=rv64gcv_zvl256b_zicbop_zicboz}.

\begin{table}[htbp]
  \centering\small
  \caption{Software prefetch in the lab (f32, one session, fastest step).}\label{r03:tab:pf}
  \begin{tabularx}{\linewidth}{@{}Xrrr@{}}
    \toprule
    \textbf{Kernel} & \textbf{MLUP/s} & \textbf{cycles/cell} & \textbf{instr./cell}\\
    \midrule
    RVV with CSE (as before this work) & 2.47 & 653 & 50.5\\
    + prefetch 64 / 96~B & 4.95 / 4.98 & 326 / 324 & 68\\
    \textbf{+ prefetch 128~B} & \textbf{5.19} & \textbf{311} & 68\\
    + prefetch 256 / 512~B & 4.85 / 4.03 & 332 / 399 & 68\\
    + prefetch 128~B, reads only & 4.91 & 333 & 58\\
    scalar with CSE & 1.57 & 1035 & 311\\
    scalar \emph{without} CSE + prefetch 128~B (one prefetch per 16 cells)$^{a}$ & 1.81 & 887 & 380\\
    RVV with CSE, density/velocity stores removed (not a complete kernel) & 3.19 & 559 & 48\\
    \bottomrule
  \end{tabularx}

  \smallskip
  \begin{minipage}{0.92\linewidth}\footnotesize
    $^{a}$ Named \code{f32-generic-cse_scalar+pf=128} in the lab and first recorded as ``scalar
    with CSE + prefetch'', but generated without global CSE: a naming error in the lab generator,
    found while writing this report (\cref{r03:sec:repro-gen}).
  \end{minipage}
\end{table}

Prefetch alone doubles the vector kernel. 128~bytes ahead is the best distance, and it stayed
the best in every later measurement.

For the scalar kernel the lab has no clean comparison. The scalar kernel with prefetch was, by
mistake, generated without CSE (note $a$). It ran 15\,\% faster than the scalar kernel with CSE
and without prefetch (1.81 against 1.57~MLUP/s), although it executes more instructions per cell
(380 against 311). So prefetch helps a scalar kernel too, but its gain on the scalar kernel
\emph{with} CSE was not measured. The application cannot use it in any case: the \ps{} option
exists for RVV kernels only (\cref{r03:sec:ps-prefetch}).

The last row removes the four stores of density and velocity. The PDF update is unchanged,
but the kernel no longer writes the output, so it is not a complete kernel for the benchmark,
which needs the output once at the end. It saves 14\,\% of the cycles. This led to the kernels
without output of \cref{r03:sec:wl-sweepgen}, which are used in all steps except the last.

\section{Fewer streams: the direction split}\label{r03:sec:lever-split}

\ghws{misc/x60-lab/lab/hand/f32-hand-split.cpp} computes the same update as the generated
kernel, but in three kinds of pass over a small block of rows instead of one pass over all 19
directions. The block is BY rows of one $z$-plane (BY is a parameter). For each block:
\begin{enumerate}
  \item \textbf{moment passes}: each pass reads the PDFs of a group of GM directions and adds
        their contributions to density and momentum into a small scratch array, which stays in
        the L1 cache;
  \item \textbf{macroscopic pass}: computes $\rho$ and $u$ from the sums and writes the
        output;
  \item \textbf{collision passes}: each pass reads a group of directions again, which now hit
        in L1 or L2, collides them and stores the results.
\end{enumerate}
Each pass has only a few live streams, which the hardware prefetcher can follow. With
BY~=~2 and groups of 5 directions the hand kernel reaches 6.98~MLUP/s (fastest step; 237
cycles per cell), against 5.2 for the generated kernel with prefetch. L2 misses caused by
loads drop from 2.45 to 0.70 per cell (hardware counters). Software prefetch makes the split
kernel slower, because it competes with the hardware prefetcher. Blocks of 8 rows no longer
fit in L1: 3.85~MLUP/s (measured with prefetch).

\section{In-place AA streaming}\label{r03:sec:lever-aa}

AA streaming (\cref{r03:sec:bg-stream}) reads and writes the same array. That removes the 19
write streams into the second array: 19 + 4 streams remain instead of 19 + 23
(\cref{r03:fig:streams-schemes}). It also moves fewer bytes, about 184 instead of about 260
per cell in f32, because there is no write-allocate of a second array.

\begin{figure}[htbp]
  \centering
  \resizebox{\linewidth}{!}{% Memory streams of one cell update in each sweep variant, against the ~8 streams the X60's hardware
% prefetcher follows. One square = one stream. Bytes per cell for f32.
\begin{tikzpicture}[every node/.style={font=\scriptsize},
    st/.style={minimum size=2.4mm,inner sep=0pt,draw=white,line width=0.3pt},
    R/.style={st,fill=psblue!70},
    W/.style={st,fill=psorange!80},
    RW/.style={st,fill=psteal!75},
    O/.style={st,fill=psgray!55}]
  \def\rthreeW{0.28}
  % row: label, then squares
  \newcommand{\rthreerow}[4]{% y, label, list of style:count, total text
    \node[anchor=east,align=right,text width=34mm] at (-0.15,#1) {#2};
    \xdef\rthreeN{0}%
    \foreach \sty/\cnt in {#3} {
      \foreach \k in {1,...,\cnt} {
        \pgfmathsetmacro{\xx}{\rthreeN*\rthreeW+\rthreeW/2}
        \node[\sty] at (\xx,#1) {};
        \pgfmathtruncatemacro{\nn}{\rthreeN+1}\xdef\rthreeN{\nn}%
      }
    }
    \node[anchor=west] at (\rthreeN*\rthreeW+0.1,#1) {#4};
  }
  \rthreerow{0}{pull, $\approx$260~B/cell}{R/19,W/19,O/4}{42}
  \rthreerow{-0.55}{AA, $\approx$184~B/cell}{RW/19,O/4}{23}
  \rthreerow{-1.10}{AA without $\rho$/$u$ output,\\152~B/cell}{RW/19}{19}
  \rthreerow{-1.75}{split AA, largest pass}{RW/7}{7 (+ L1 scratch)}
  % the prefetcher's limit
  \draw[r03landed!80!black,thick,dashed] (8*\rthreeW,0.35) -- (8*\rthreeW,-2.05)
    node[below,text=r03landed!70!black] {about 8: what the hardware prefetcher follows};
  % legend
  \begin{scope}[yshift=-2.75cm]
    \node[R] at (0.1,0) {}; \node[anchor=west] at (0.2,0) {read};
    \node[W] at (1.5,0) {}; \node[anchor=west] at (1.6,0) {write (second field)};
    \node[RW] at (5.0,0) {}; \node[anchor=west] at (5.1,0) {read and written in place};
    \node[O] at (9.3,0) {}; \node[anchor=west] at (9.4,0) {output $\rho$, $u_x$, $u_y$, $u_z$};
  \end{scope}
\end{tikzpicture}
}
  \caption[Memory streams per kernel structure]{Concurrent memory streams of one cell update
    for each kernel structure, against the about 8 that the X60's hardware prefetcher follows.
    One square is one stream. Bytes per cell for f32.}\label{r03:fig:streams-schemes}
\end{figure}

\paragraph{Hand kernel.} \ghws{misc/x60-lab/lab/hand/hand-aa.hpp} combines AA streaming with
the direction split. Its collision passes work on pairs of opposite directions $(i,
\bar\imath)$: in the AA pattern, the two places a pair writes are the two places it read, so
both values must be loaded before either is stored. GP is the number of pairs per collision
pass. In f32 with BY~=~2, GM~=~5 and GP~=~3 the kernel reaches 9.11~MLUP/s (fastest step), 180
cycles per cell. f64 runs best with one-row blocks (BY~=~1): 4.26~MLUP/s, 416 cycles per cell.
Memory is no longer the limit (0.22 L2 misses from loads per cell). The largest source of
stalls is now the issue queue of the vector floating-point unit.

\paragraph{Generated kernel.} lbmpy already implements the AA pattern
(\code{streaming_pattern="aa"}): it produces one kernel for even and one for odd steps. Generated
through the \ps{} RVV back end with prefetch 128~bytes:

\begin{table}[htbp]
  \centering\small
  \caption{lbmpy's AA kernels generated for RVV (one session, fastest step, instruction
    selection as it was before \cref{r03:sec:lever-selector}).}\label{r03:tab:genaa}
  \begin{tabular}{@{}lrrrr@{}}
    \toprule
    & \multicolumn{2}{c}{\textbf{f32}} & \multicolumn{2}{c}{\textbf{f64}}\\
    \cmidrule(lr){2-3}\cmidrule(l){4-5}
    \textbf{Kernel} & \textbf{MLUP/s} & \textbf{cycles/cell} & \textbf{MLUP/s} & \textbf{cycles/cell}\\
    \midrule
    pull + prefetch 128~B & 4.50 & 357 & 2.44 & 682\\
    AA, no prefetch & 2.81 & 651 & 1.41 & 1290\\
    \textbf{AA + prefetch 128~B} & \textbf{7.18} & \textbf{235} & \textbf{3.86} & \textbf{451}\\
    \bottomrule
  \end{tabular}
\end{table}

The generated AA kernel with prefetch is \sx{1.60} (f32) and \sx{1.58} (f64) faster than pull
with prefetch. Without prefetch, AA's 23 streams are still too many for the hardware
prefetcher. Under QEMU the generated f64 AA kernel gives bit-identical results to the generated
f64 pull kernel, which also confirms the lab driver's ghost-layer handling for AA.

\section{Register grouping and instruction scheduling}\label{r03:sec:lever-sched}

\paragraph{LMUL.} With LMUL~2, each instruction processes two registers, so the kernel needs
half as many instructions. The hand AA kernel nevertheless becomes slower: 216 instead of 163
cycles per cell. The X60 executes an LMUL-2 instruction as two consecutive halves, so a chain
of dependent instructions takes no less time per cell. And with registers grouped in pairs, only
16 groups remain, so fewer independent chains can be interleaved to hide latencies. LMUL~4 is
slower still (226--229 cycles per cell), because the kernel no longer fits in the registers.
Rejected.

\paragraph{X60 scheduling model.} Clang~18 can order instructions for the X60
(\code{-mtune=spacemit-x60}, \cref{r03:sec:bg-rvv}); GCC~14 cannot.

\begin{table}[htbp]
  \centering\small
  \caption{Compilers and the X60 scheduling model, MLUP/s (one session, fastest step).}\label{r03:tab:tune}
  \begin{tabular}{@{}lrrr@{}}
    \toprule
    \textbf{Kernel} & \textbf{GCC 14} & \textbf{Clang 18} & \textbf{Clang 18, X60 model}\\
    \midrule
    f32 scalar with CSE & 1.46 & 1.19 & 1.56\\
    f32 pull + prefetch & 5.12 & 5.55 & 5.45\\
    f32 generated AA + prefetch & \textbf{8.99} & 7.41 & 8.09\\
    f32 hand AA (BY~=~2, GP~=~3) & 10.13 & 8.83 & \textbf{11.10}\\
    f64 scalar with CSE & 1.03 & 0.97 & 1.07\\
    f64 pull + prefetch & 2.60 & 2.56 & 2.82\\
    f64 generated AA + prefetch & \textbf{4.11} & 3.34 & 3.94\\
    f64 hand AA (BY~=~1, GP~=~3) & 4.27 & 3.81 & \textbf{4.63}\\
    \bottomrule
  \end{tabular}
\end{table}

The scheduling model makes every Clang build faster, by 9 to 31\,\%, the scalar baseline
included. For the generated AA kernels GCC~14 remains faster than Clang with the model, so the
end-to-end results in \cref{r03:ch:results} are GCC.

\section{Instruction selection: .vf operands and FMA}\label{r03:sec:lever-selector}

\term{Instruction selection} is the step of \ps{} that chooses which RVV intrinsic implements
each operation. With the original selection, the generated AA loop had 471 instructions per 8
cells, of which only 280 were vector instructions (inspected with the GCC~16 cross compiler).
It used no FMA at all, and reloaded 23 vector registers from the stack and saved 7 to it
(\term{spills}) in every iteration. The reason: the 31 constants of the update were kept as
vector registers and crowded out the 32 vector registers. \Cref{r03:sec:ps-selector}
describes the change; its effect on the board (log \file{lab3.log}):

\begin{table}[htbp]
  \centering\small
  \caption{Original $\rightarrow$ new instruction selection (one session; MLUP/s: fastest step,
    cycles per cell: average over all steps).}\label{r03:tab:selector}
  \begin{tabular}{@{}lrrr@{}}
    \toprule
    \textbf{Kernel} & \textbf{MLUP/s} & \textbf{cycles/cell} & \textbf{instr./cell}\\
    \midrule
    f32 pull + prefetch & 4.85 $\rightarrow$ 4.82 & 334 $\rightarrow$ 337 & 68.0 $\rightarrow$ 64.2\\
    f64 pull + prefetch & 2.44 $\rightarrow$ 2.59 & 685 $\rightarrow$ 642 & 131.8 $\rightarrow$ 124.2\\
    f32 AA + prefetch & 7.31 $\rightarrow$ \textbf{8.88} & 234 $\rightarrow$ \textbf{192} & 70.5 $\rightarrow$ 51.5\\
    f64 AA + prefetch & 3.86 $\rightarrow$ 4.09 & 461 $\rightarrow$ 433 & 138.2 $\rightarrow$ 100.8\\
    \bottomrule
  \end{tabular}
\end{table}

The change pays off where the kernel is no longer waiting for memory (AA with prefetch). In the
pull kernels the extra instructions had been hidden behind memory stalls.

\section{Row pointers and constants in registers}\label{r03:sec:lever-rp}

After the new instruction selection, two inefficiencies of the compiled code remained in the
AA loop (GCC~16 cross compiler):
\begin{itemize}
  \item An access to the neighbour $x \pm 1$ has an address offset of $\pm 4$~bytes (f32). RVV
        loads and stores cannot add an offset to an address register, so GCC computes such
        addresses from the loop counter with an extra addition in every iteration, and spills
        some of them.
  \item GCC loads every scalar operand of a .vf instruction from the constant table in memory
        in every iteration (\code{lla} + \code{flw}), instead of keeping it in one of the
        floating-point registers, which a vector loop otherwise leaves unused.
\end{itemize}
The fix computes a pointer to the start of each row once before the $x$ loop (\term{row
pointers}), and hides the constants' values from the compiler so that it keeps them in
registers (\cref{r03:sec:ps-rowptr}). It was tried first as an edit of the generated source,
then implemented in \ps{}. It reduces the f32 odd AA step from 383 to 305 instructions per 8
cells.

On the board, with GCC~14, the lab kernels execute 1.7 (f32) to 3.7 (f64) fewer instructions
per cell, but their cycles do not drop: f32 AA $-0.4$\,\%, f64 AA $-2$\,\%, f32 pull
$+1.1$\,\% (log \file{lab5.log}, mean of 2 rounds, which differ by up to 2\,\%). In the lab
the kernels wait for memory, not for integer instructions.

The application disagrees in part. There the variants with prefetch became faster
(\cref{r03:tab:res-cg}): pull with prefetch by 5.2\,\% (f32) and 8.7\,\% (f64), AA with
prefetch by 2.5\,\% and 2.2\,\%. One difference between the two measurements: the lab used the
hand-edited source, the application the code generated by \ps{}. The change was kept
(\cref{r03:sec:res-cg}).

\section{Alignment, huge pages and padding}\label{r03:sec:lever-layout}

\begin{table}[htbp]
  \centering\small
  \caption{Interior and rows aligned to 16~bytes (waLBerla's layout), 32 or 64~bytes, and
    32~bytes with transparent huge pages. MLUP/s over all steps (\texttt{mean\_mlups}), mean of 2
    rounds in one session (log \texttt{lab5.log}).}\label{r03:tab:align}
  \begin{tabular}{@{}lrrrr@{}}
    \toprule
    \textbf{Kernel} & \textbf{16~B} & \textbf{32~B} & \textbf{64~B} & \textbf{32~B + THP}\\
    \midrule
    f32 scalar (pull) & 1.62 & 1.62 & 1.69 & 1.57\\
    f32 pull + prefetch & 5.16 & 5.25 & 3.67 & 3.93\\
    f32 AA + prefetch & \textbf{8.92} & 8.34 & 6.81 & 7.55\\
    f64 scalar (pull) & 0.885 & 1.01 & 1.185 & 1.02\\
    f64 pull + prefetch & 2.28 & 2.46 & 2.04 & 2.48\\
    f64 AA + prefetch & \textbf{3.74} & 3.65 & 3.10 & 3.63\\
    \bottomrule
  \end{tabular}
\end{table}

waLBerla's 16-byte layout is the best one for the AA kernels; 64-byte alignment costs the vector
kernels 11 to 29\,\%. The scalar pull kernels gain from 64-byte alignment (f64 +34\,\%, f32
+4\,\%), but the best scalar build is an AA kernel, which this does not affect. Rejected.

With huge pages (and 32-byte alignment), the f32 AA odd step is 14\,\% faster and the even step
24\,\% slower than with 32-byte alignment and 4~KiB pages. A possible explanation is conflicts
between the 19 PDF slices in the cache sets or the DRAM banks, once 2~MiB pages fix their
physical offsets to each other. The test of that explanation, extra padding between the
slices, is pending (\cref{r03:ch:open}).

\section{The direction split, generated from lbmpy}\label{r03:sec:lever-gensplit}

The hand split kernels showed what the structure can do, but they are written by hand, for one
method. \Cref{r03:sec:wl-split} generates the same structure from lbmpy. Under QEMU the result is
correct: for every BY, GM and GP tried, the f64 split kernels agree with the generated AA kernel
in all printed digits. Their inner loops have 418 instructions per 8 cells over all 8 passes
(even step, f32, BY~=~2, GM~=~5, GP~=~3; GCC~16 cross compiler), against 340 for the generated
single-pass AA kernel and 469 for the hand kernel. Whether fewer live streams outweigh the extra
instructions on the board, as they did for the hand kernel, was still being measured
(\cref{r03:ch:open}).
